\subsection{序域上的向量空间}

设 K 为一个序域，E 为 K 上的一个向量空间。关系«存在 \(\lambda > 0\) 在 K 中使得 \(y = \lambda x \gg\) ，在集合 \(E - \{0\}\) 的两个元素 x 与 y 之间，是一个等价关系。在此关系下的等价类称为以 0 为原点的开半直线；一个开半直线与 \(\{0\}\) 的并称为以 0 为原点的闭半直线（有时简称半直线）。每个向量 \(a \neq 0\) ，若它包含在一条开（相应地，闭）半直线 \(\Delta\) 中，则称为 \(\Delta\) 的方向向量，而 \(\Delta\) 是向量 \(\lambda a\) （对所有标量 \(\lambda > 0\) ，相应地 \(\lambda \geqslant 0\) ）所成的集合。每条过 0 的直线 D 恰好包含两条以 0 为原点的开（相应地，闭）半直线；如果 \(\Delta\) 是其中之一，则 \(-\Delta\) 是另一个（称为 \(\Delta\) 的反向半直线）。

现在如果 F 是 K 上的一个仿射空间，且 E 是 F 的平移空间，那么 F 的任何形如 \(\Delta = a + \Delta_{0}\) 的子集，其中 \(\Delta_{0}\) 是 E 的一条开（相应地，闭）半直线，称为以原点 \(a \in F\) 为起点的开（相应地，闭）半直线。半直线 \(\Delta_{0}\) 完全由 \(\Delta\) 确定（因为它是对任意 \(b \neq a\) 属于 \(\Delta\) 而言方向向量为 b - a 的那条半直线），称为 \(\Delta\) 的方向；\(\Delta_{0}\) 的方向向量也称为 \(\Delta\) 的方向向量。

现在假设 E 在 K 上具有有限维数 n；那么已知（III，第 518 页，推论 1）第 n 次外幂 \(\bigwedge^{n}E\) 是 K 上的一维向量空间，因此是以 0 为原点的两条相反闭半直线之并。这些半直线称为 E 的定向；空间 E 连同给定的其中一条半直线 \(\Delta\) 称为已定向的；一个 \(n\) -向量 \(z\) 在此定向下，若它属于 \(\Delta\) （相应地，属于 \(-\Delta\) ），则称为正的（相应地，负的）；在相反的定向下，它是负的（相应地，正的）。

仿射空间 F 在 K 上的一个定向，按定义是 F 的平移空间的一个定向；空间 F 连同这样的定向，称为有向仿射空间。

设 \(\mathbf{E}\) 是 \(\mathbf{K}\) 上一个已定向的向量空间，维数为 \(n\) ；有序基 \((a_i)_{1 \leqslant i \leqslant n}\) （\(\mathbf{E}\) 的）称为正的或顺向的（相应地，负的或逆向的），如果 \(n\) -向量 \(a_1 \wedge a_2 \wedge \ldots \wedge a_n\) 为正（相应地，为负）。若 \(u\) 是向量空间 \(\mathbf{E}\) 的一个自同构，则 \((\bigwedge^{n} u)(z) = \det(u) \cdot z\) 对所有 \(z \in \bigwedge^{n} \mathbf{E}\) 成立，故

\(\bigwedge^{n} u\) 保持 E 的定向不变（或者，如我们所说，保持定向）的一个充分必要条件是 \(\det(u) > 0\) ；具有此性质的自同构恰好是 E 作为已定向向量空间的自同构；它们构成一个正规子群 \(\mathbf{GL}^{+}(\mathbf{E})\) ，即线性群 \(\mathbf{GL}(\mathbf{E})\) 的子群，只要 \(E \neq 0\) ，其指数为 2。

当 \(\mathbf{E} = 0\) 时，\(\mathbf{GL}(\mathbf{E}) = \operatorname{End}(\mathbf{E})\) 仅包含恒等映射 \(1_{\mathbf{E}}\) ，且按定义 \(\det(1_{\mathbf{E}}) = 1\) 。注意，在此情形下按定义 \(\bigwedge^{n}\mathbf{E} = \bigwedge^{0}\mathbf{E} = \mathbf{K}\) ；\(\mathbf{K}\) 的由正标量组成的半直线称为零空间的标准定向。

设 \(\mathbf{M}\) 与 \(\mathbf{N}\) 是维数分别为 \(p\) 与 \(n - p\) 的两个互补子空间，属于向量空间 \(\mathbf{E}\) （维数为 \(n\) ）；如果 \(z'\) （相应地 \(z''\) ）是 \(\bigwedge^{p} \mathbf{M}\) （相应地 \(\bigwedge^{n - p} \mathbf{N}\) ）中的一个非零向量，则 \(z' \wedge z''\) 是

\(\bigwedge^{n} \mathbf{E}\) 中的非零向量。给定 M 上的一个定向和 N 上的一个定向，向量 \(z' \wedge z''\) 对正的 \(z'\) 与 \(z''\) 构成 E 的一个定向，称为 M 的定向与 N 的定向的乘积定向（当 \(p(n - p)\) 为奇数时，它依赖于因子的顺序）。反之，给定 E 和 M 上的定向，存在 N 上唯一的定向，使得 E 上给定的定向是 M 上给定的定向与该 N 上的定向（按此顺序）的乘积；该定向称为关于 E 的定向的 M 的定向的互补定向。若 \(N'\) 是 M 的第二个互补子空间，则平行于 M 的典范投影 \(N \to N'\) 将 N 的互补定向映到 \(N'\) 的互补定向。因此，N 的互补定向在典范映射 \(N \to E/M\) 下的像不依赖于补空间 N 的选择；它称为 E 的定向除以 M 的定向所得的 E/M 上的商定向。

